% !TeX root = ../../Book4.tex
% 纲要标题：### 19.4 扭对与 Happel–Reiten–Smalø 倾斜
% 纲要位置：计划纲要.md，第 2999 行。

\section{扭对与 Happel–Reiten–Smalø 倾斜}
\label{b4:sec:1904}


% 纲要内容（仅作写作计划，尚非教材正文）：
%
% * 在交换范畴 \(\mathcal A\) 中定义扭对 \((\mathcal T,\mathcal F)\)：\(\operatorname{Hom}(\mathcal T,\mathcal F)=0\)，每个对象有扭部分与无扭商组成的短正合列
% * 在 \(D^b(\mathcal A)\) 中由标准截断和扭分解构造倾斜后的两个子范畴，逐项证明正交性与截断三角公理
% * 固定本书方向：新心由仅在 \(-1,0\) 次有上同调、且 \(H^{-1}\in\mathcal F\)、\(H^0\in\mathcal T\) 的对象组成，即 \(\langle\mathcal F[1],\mathcal T\rangle\)
% * 利用第19.3节已证的心的交换范畴结构解释新心的核、余核和扩张；区分扭对倾斜与第20章倾斜对象的不同数据
% * 新的心不自动给出与原交换范畴的导出等价；这需要额外条件，参见 [V13-6]
%

\subsection{扭对与扭分解}
\label{b4:subsec:1904:01}

标准心把每个模或交换范畴对象放在零次。若某一类对象适合放在相邻次数，问题便成为：怎样移动这一类而仍保留核、余核和短正合列？扭分解提供了与任意对象都相容的分类方法。

\begin{definition}\label{b4:def:torsion-pair}
设 \(\mathcal A\) 为交换范畴，\(\mathcal T,\mathcal F\) 为在同构下封闭的全子范畴。若 \(\operatorname{Hom}_{\mathcal A}(T,F)=0\) 对 \(T\in\mathcal T,F\in\mathcal F\) 成立，且每个 \(M\in\mathcal A\) 都有短正合列
\[
 0\longrightarrow tM\longrightarrow M\longrightarrow fM\longrightarrow0,
 \qquad tM\in\mathcal T,\quad fM\in\mathcal F,
\]
则称 \((\mathcal T,\mathcal F)\) 为 \(\mathcal A\) 中的\term[niudui]{扭对}[torsion pair]，此列称为扭分解。这里的扭与无扭是对这对子范畴的命名，不默认它们由某个环的非零元是否零化元素来定义。
\end{definition}

\begin{proposition}\label{b4:prop:torsion-pair-properties}
设 \((\mathcal T,\mathcal F)\) 为交换范畴中的扭对。扭分解唯一到保持 \(M\) 的唯一同构，并且对 \(M\) 函子性变化；\(\mathcal T={}^\perp\mathcal F\)、\(\mathcal F=\mathcal T^\perp\)。此外，\(\mathcal T\) 对商及扩张封闭，\(\mathcal F\) 对子对象及扩张封闭。
\end{proposition}
\begin{proof}
对任意 \(T\in\mathcal T\)，映射 \(T\to M\to fM\) 为零，所以唯一经 \(tM\) 分解；这使 \(tM\to M\) 成为所有此类映射的泛对象。对 \(M\to F\)、\(F\in\mathcal F\)，映射在 \(tM\) 上为零，所以唯一经 \(fM\) 分解。这证明唯一性与函子性。若 \(M\) 与所有 \(F\) 正交，则其满射 \(M\to fM\) 为零，故 \(fM=0\) 且 \(M\in\mathcal T\)；另一侧对偶。商封闭性由 \(\operatorname{Hom}(Q,F)\to\operatorname{Hom}(T,F)\) 单射得到，子对象封闭性对偶。若两个端点都属于同一正交类，短正合列的 Hom 正合性使中间项也属于该类，故有扩张封闭性。
\end{proof}

交换群中的挠群与无挠群给出一个扭对：\(tM\) 是所有有限阶元素组成的子群，\(M/tM\) 无挠，因为 \(n\bar x=0\) 意味着 \(nx\) 有限阶，进而 \(x\) 有限阶。下一节的扭对则由线性映射的两个端点定义，与整数挠元没有关系。

\subsection{扭对倾斜的构造与公理}
\label{b4:subsec:1904:02}

\begin{theorem}\label{b4:thm:hrs-tilt}
设 \(\mathcal A\) 为交换范畴，\((\mathcal T,\mathcal F)\) 是其中的扭对。在 \(D^b(\mathcal A)\) 中令
\begin{align*}
 \mathcal D_t^{\le0}&=\{X:H^i(X)=0\ (i>0),\ H^0(X)\in\mathcal T\},\\
 \mathcal D_t^{\ge0}&=\{X:H^i(X)=0\ (i<-1),\ H^{-1}(X)\in\mathcal F\}.
\end{align*}
则 \((\mathcal D_t^{\le0},\mathcal D_t^{\ge0})\) 是有界 t-结构。这个构造称为\term[hrsqingxie]{Happel--Reiten--Smalø 倾斜}[Happel--Reiten--Smalø tilt]，简称 HRS 倾斜。
\end{theorem}
\begin{proof}
先检查移位条件。若 \(X\in\mathcal D_t^{\le0}\)，则 \(X[1]\) 的正次上同调与零次上同调均为零，所以仍在此类；若 \(Y\in\mathcal D_t^{\ge0}\)，则 \(Y[-1]\) 在低于零次的上同调为零，故仍在此类。

新结构的 \(\mathcal D_t^{\ge1}=\mathcal D_t^{\ge0}[-1]\) 由满足 \(H^i(Y)=0\ (i<0)\)、\(H^0(Y)\in\mathcal F\) 的对象组成。若 \(X\in\mathcal D_t^{\le0}\)、\(Y\in\mathcal D_t^{\ge1}\)，标准截断的伴随首先把任意 \(X\to Y\) 唯一分解为 \(X\to H^0(X)\to Y\)，再把后一个映射唯一分解为 \(H^0(X)\to H^0(Y)\to Y\)。中间映射从 \(\mathcal T\) 到 \(\mathcal F\)，所以为零。这证明新正交性。

为构造新截断，给定 \(X\)，记 \(A=\tau^{\le0}X\)、\(T=tH^0(X)\)、\(F=fH^0(X)\)。取复合 \(A\to H^0(X)\to F\)，完成为三角
\[
 U\longrightarrow A\longrightarrow F\longrightarrow U[1].
\]
普通上同调长正合列表明 \(H^i(U)=H^i(X)\) 当 \(i<0\)，\(H^0(U)=T\)，\(H^i(U)=0\) 当 \(i>0\)；这里用到了 \(H^0(X)\to F\) 是满射。故 \(U\in\mathcal D_t^{\le0}\)。把 \(U\to A\to X\) 的复合完成为三角 \(U\to X\to V\)，八面体公理另给出三角
\[
 F\longrightarrow V\longrightarrow\tau^{\ge1}X\longrightarrow F[1].
\]
由它得 \(H^i(V)=0\) 当 \(i<0\)，且 \(H^0(V)=F\)，所以 \(V\in\mathcal D_t^{\ge1}\)。分解公理因此成立。

最后有包含关系 \(D^{\le-1}\subseteq\mathcal D_t^{\le0}\subseteq D^{\le0}\) 及 \(D^{\ge0}\subseteq\mathcal D_t^{\ge0}\subseteq D^{\ge-1}\)。标准 t-结构在有界导出范畴中有界，故新结构也有界。
\end{proof}

\subsection{倾斜方向与新心的上同调条件}
\label{b4:subsec:1904:03}

\begin{corollary}\label{b4:cor:hrs-heart-description}
设 \(\mathcal A\) 为交换范畴，\((\mathcal T,\mathcal F)\) 为扭对，\(\mathcal B\subseteq D^b(\mathcal A)\) 为\cref{b4:thm:hrs-tilt}给出的新心。则 \(\mathcal B\) 由恰好满足
\[
 H^i(E)=0\quad(i\ne-1,0),\qquad
 H^{-1}(E)\in\mathcal F,\quad H^0(E)\in\mathcal T
\]
的对象组成。每个 \(E\in\mathcal B\) 都有新心中的短正合列
\[
 0\longrightarrow H^{-1}(E)[1]\longrightarrow E
 \longrightarrow H^0(E)\longrightarrow0.
\]
因此常记 \(\mathcal B=\langle\mathcal F[1],\mathcal T\rangle\)，尖括号在这里表示按上述次序取扩张。
\end{corollary}
\begin{proof}
两个新子范畴的交给出列出的上同调条件。标准截断三角为 \(H^{-1}(E)[1]\to E\to H^0(E)\to H^{-1}(E)[2]\)，三项都在新心中，由\cref{b4:prop:heart-exact-triangles}即成短正合列。反之，以 \(F[1]\)、\(T\) 为两端的任何三角扩张都在新心中，因为心对扩张封闭。
\end{proof}

这里移动的是 \(\mathcal F\)：它的对象由零次移到 \(-1\) 次。文献也常采用反向约定，得到 \(\langle\mathcal F,\mathcal T[-1]\rangle=\mathcal B[-1]\)。比较结论时须先换算方向，不能只看“倾斜心”这一名称。

\subsection{新心的核、余核与扩张}
\label{b4:subsec:1904:04}

新心已经是交换范畴；它的核与余核由\cref{b4:thm:heart-abelian}计算，但那里要采用新的 \(H_t^i\)。普通上同调 \(H^i\) 只用来说明新心的对象是什么，并不直接给出新心中每个态射的核与余核。

\begin{example}\label{b4:exm:hrs-integer-multiplication}
取 \(\mathcal A=\mathbb Z\text{-}\mathrm{Mod}\)，扭对为挠群与无挠群，\(p\) 为素数。对象 \(\mathbb Z[1]\) 位于新心。普通模上的单射 \(p:\mathbb Z\to\mathbb Z\)，在新心中成为满态射
\[
 p:\mathbb Z[1]\longrightarrow\mathbb Z[1],\qquad
 \ker_{\mathcal B}p=\mathbb Z/p,
\]
并给出短正合列
\[
 0\longrightarrow\mathbb Z/p\longrightarrow\mathbb Z[1]
 \xrightarrow{p}\mathbb Z[1]\longrightarrow0.
\]
事实上，标准短正合列 \(0\to\mathbb Z\xrightarrow{p}\mathbb Z\to\mathbb Z/p\to0\) 的三角旋转后，前三项就是这三个新心对象；必要时把一个结构映射改乘 \(-1\)，即与旋转符号一致。由\cref{b4:prop:heart-exact-triangles}得到所述短列。
\end{example}

\((\mathcal F[1],\mathcal T)\) 本身还是新心中的扭对。其分解由\cref{b4:cor:hrs-heart-description}给出，正交性来自 \(\operatorname{Hom}(F[1],T)=\operatorname{Hom}(F,T[-1])=0\)。这一事实说明新心并非没有原扭分解的痕迹；变化的是两类对象所在的次数，以及它们之间的正合列。

\subsection{倾斜心与导出等价的条件}
\label{b4:subsec:1904:05}

HRS 倾斜的数据是一整个交换范畴中的两类对象及其扭分解。\chapref{b4:ch:20}讨论的倾斜对象则是一项对象及其自扩张消失、生成等条件。两种构造有关联，却不能仅凭名称相同便认作同一个定义。

\ProseParagraphBreak
尤其是，在 \(D^b(\mathcal A)\) 中得到新心 \(\mathcal B\)，并不已经得到 \(D^b(\mathcal B)\simeq D^b(\mathcal A)\)。若存在与心的包含相容的三角等价，它对 \(X,Y\in\mathcal B\) 必须给出每个次数的同构 \(\operatorname{Hom}_{D^b(\mathcal B)}(X,Y[n])\simeq\operatorname{Hom}_{D^b(\mathcal A)}(X,Y[n])\)。当新心有足够投射或内射对象时，左边可用其中的消解计算为 Ext。一次扩张与环境态射的对应已由\cref{b4:prop:heart-ext1}保证，高阶的一致性却需要另作证明。

\ProseParagraphBreak
下一节直接计算一个两顶点范畴：新心为半单范畴，故心内二次 Ext 为零，而环境的某个二次移位 Hom 为 \(k\)。这将排除任何与该包含相容的导出等价，而不把“倾斜”误当作已有等价的凭据。本章不借此陈述未经证明的一般倾斜等价判据。
